\subsection{结合律}

定义5. 合成律 \((x,y)\mapsto x\top y\) （定义在集合 \(\mathbf{E}\) 上）称为结合的，如果对于所有元素 \(x,y,z\) （在 \(\mathbf{E}\) 中），

\[
(x \top y) \top z = x \top (y \top z).
\]

合成律为结合律的广群称为结合广群。

结合律的反合成律也是结合的。

例子。(1) 自然数的加法和乘法是 \(\mathbf{N}\) 上的结合合成律（集合论，III，§ 3，第 3 号，命题 5 的推论）

\begin{enumerate}
\def\labelenumi{(\arabic{enumi})}
\setcounter{enumi}{1}
\tightlist
\item
  第1号中例(1)、(3)和(4)所引用的律都是结合律。
\end{enumerate}

定理1（结合性定理）。设 E 为一个结合广群，其律记为 \(\top\) 。设 A 是一个全序非空有限集，它是非空子集的一个有序序列 \((\mathbf{B}_l)_{i \in I}\) 的并集，使得关系 \(\alpha \in \mathbf{B}_i\) , \(\beta \in \mathbf{B}_j\) , \(i < j\) 蕴含 \(\alpha < \beta\) ；设 \((x_\alpha)_{\alpha \in A}\) 为以 A 为指标集的 E 中元素的一个有序序列。则

\[
\bigcap_ {\alpha \in A} x _ {\alpha} = \bigcap_ {i \in I} \left(\bigcap_ {\alpha \in B _ {i}} x _ {\alpha}\right)\tag{3}
\]

† 使用这一术语及记号 \(\prod_{\alpha \in \mathbf{A}} x_{\alpha}\) ，若存在与在集合论（集合论，II，§ 5，第 3 号）中定义的诸集合 \(x_{\alpha}\) 的乘积混淆的风险，则必须避免。然而，如果 \(x_{\alpha}\) 是基数且加法（相应地，乘法）是基数和（相应地，基数积），则在上述记号中记作 \(\sum_{\alpha \in \mathbf{A}} x_{\alpha}\) （相应地 \(\prod_{\alpha \in \mathbf{A}} x_{\alpha}\) ）的基数，是族 \((x_{\alpha})_{\alpha \in \mathbf{A}}\) 的基数和（相应地，基数积）（集合论，III，§ 3，第 3 号）。

我们通过对 A 的基数 n 进行归纳来证明该定理。设 p 为 I 的基数，h 为其最小元素；设 \(J = I - \{h\}\) 。若 n = 1，由于 \(B_{i}\) 非空，必然有 p = 1，定理显然成立。否则，假设定理对至多含 n - 1 个元素的指标集成立，我们区分两种情况：

\begin{enumerate}
\def\labelenumi{(\alph{enumi})}
\tightlist
\item
  \(B_{h}\) 只有一个元素 \(\beta\) 。设 \(C = \bigcup_{i \in J} B_{i}\) 。(3)的左边根据定义等于 \(x_{\beta} \top \left( \bigcap_{\alpha \in C} x_{\alpha} \right)\) ；右边根据定义等于
\end{enumerate}

\[
x _ {\beta} \top \left(\underset {i \in J} {\mathsf {T}} \left(\underset {\alpha \in B _ {i}} {\mathsf {T}} x _ {\alpha}\right)\right);
\]

结果由定理对 C 和 \((\mathbf{B}_{i})_{i\in J}\) 成立这一假设推出.

\begin{enumerate}
\def\labelenumi{(\alph{enumi})}
\setcounter{enumi}{1}
\tightlist
\item
  否则，设 \(\beta\) 为 A 的最小元素（因此也是 \(\mathbf{B}_h\) 的最小元素）；设 \(\mathbf{A}' = \mathbf{A} - \{\beta\}\) ，并设 \(\mathbf{B}_i' = \mathbf{A}' \cap \mathbf{B}_i\) （对于 \(i \in \mathbf{I}\) ）；于是 \(\mathbf{B}_i' = \mathbf{B}_i\) （对于 \(i \in \mathbf{J}\) ）。集合 \(\mathbf{A}'\) 有 \(n - 1\) 个元素，且定理的条件对 \(\mathbf{A}'\) 及其子集 \(\mathbf{B}_i'\) 成立；因此根据归纳假设：
\end{enumerate}

\[
\underset {\alpha \in A ^ {\prime}} {T} x _ {\alpha} = \left(\underset {\alpha \in B ^ {\prime} h} {T} x _ {\alpha}\right) \top \left(\underset {i \in J} {T} \left(\underset {\alpha \in B _ {i}} {T} x _ {\alpha}\right)\right).
\]

将 \(x_{\beta}\) 与两边作合成，则左边根据定义有 \(\bigcap_{\alpha \in A} x_{\alpha}\) ，而右边利用结合律，得

\[
\left(x _ {\beta} \top \left(\underset {\alpha \in B ^ {\prime} h} {\top} x _ {\alpha}\right)\right) \top \left(\underset {i \in J} {\top} \left(\underset {\alpha \in B _ {i}} {\top} x _ {\alpha}\right)\right)
\]

根据定义3，这等于公式(3)的右侧。

对于记为 \(\top\) 的结合律，合成 \(\bigcap_{p\leqslant i\leqslant q}x_i\) （序列 \((x_i)_{i\in [p,q]}\) 的合成）也记作（因为不会引起混淆）

\[
x _ {p} \top \dots \top x _ {q}.
\]

定理1的一个特例是公式

\[
x _ {0} \top x _ {1} \top \dots \top x _ {n} = (x _ {0} \top x _ {1} \top \dots \top x _ {n - 1}) \top x _ {n}.
\]

考虑一个由 n 项组成的有序序列，其所有项都等于同一个元素 \(x \in E\) 。该序列的合成记作 \(\top^{n} x\) （对于记为 \(\top, \bot^{n} x\) 的律，对于记为 \(\bot\) 的律）。对于以乘法形式书写的律，该合成记作 \(x^{n}\) ，称为 x 的 n 次幂。对于以加法记号书写的律，该合成通常记作 nx。将结合性定理应用于所有项均相等的有序序列，得到等式

\[
\begin{array}{c} n _ {1} + n _ {2} + \dots + n _ {p} \\ \top \end{array} x = \binom {n _ {1}} {\top} x) \top \binom {n _ {2}} {\top} x) \top \dots \top \binom {n _ {p}} {\top} x).
\]

特别是，如果 \(p = 2\) ,

\[
{ } ^ { m + n } \top x = \left( \begin{array} { c } m \\ \top x \end{array} \right) \top \left( \begin{array} { c } n \\ \top x \end{array} \right)\tag{4}
\]

的直和，且若 \(n_1 = n_2 = \cdots = n_p = m,\)

\[
\frac {p m}{T} x = \frac {p}{T} \left(\frac {m}{T} x\right).\tag{5}
\]

如果 X 是 E 的子集，我们有时按照上述记号，用 \(\overset{p}{\top}X\) 表示集合 \(X_{1}\top X_{2}\top\cdots\top X_{p}\) ，其中

\[
\mathbf {X} _ {1} = \mathbf {X} _ {2} = \dots = \mathbf {X} _ {p} = \mathbf {X};
\]

它即所有合成 \(x_{1} \top x_{2} \top \cdots \top x_{p}\) （其中 \(x_{1} \in X, x_{2} \in X, \ldots, x_{p} \in X\) ）所成的集合.

重要的是不要将此集合与诸 \(\top x\) （其中 x 遍历 X）所成的集合混淆。

